% !TeX root = ../main.tex
\section{Proof of Bestvina-Brady Theorem}
\label{sec:morse-theory}

This section is dedicated to the proof of \cref{thm:Bestvina-Brady-finiteness-of-flag-RAAG}. 

\BestvinaBrady*
\begin{proof}
    Let \(X_L\) be the Salvetti complex associated to \(L\).
    Realize \(\phi_L\) as \(f_*\) where \(f:X_L \to \mathbb{S}^1\) is a continuous map. Let \(\tilde f:\tilde X_L \to \R\) be the map lifted to the universal cover where vertices are mapped to integers. Let base point \(x_0\) be mapped to \(0\). Then \(K_L = \tilde f ^{-1}(0)\cap \tilde X_L^{(0)}\).
    
    Notice that \(\tilde f^{-1}(0)\) is a cellular complex on which \(K_L\) acts freely and cocompactly (\(K_L\) may not be a cube complex, nor a sub-complex of \(\tilde X_L\)).

    Let \(J\subseteq \R\) and let \(Z_J:= \tilde f^{-1}(J) \subseteq \tilde X_L\).

    \begin{claim}
        \label{claim:BB-Zj-homptic-equiv}
        If \(J\subseteq J'\), both \(J, J'\) are connected, and \((J'\setminus J)\cap  \mathbb{Z}  = \emptyset\), then \(Z_J \hookrightarrow Z_{J'}\) is a homotopy equivalence.
    \end{claim}
    This is straightforward when we examine a single cell.



    \begin{definition}
        \label{ascending-descending-link}
        
    \end{definition}

\end{proof}